
\section*{Introduction}
As mentioned in the introduction to this thesis, condensed sets will be modeled by sheaves on the site $\proFiniteSets$ of profinite sets. To understand what this means we will introduce the basics of the very general theory of sites and sheaves in section \ref{section: sites and sheaves}. This will lead us to investigate general categories of sheaves, so-called \emph{Grothendieck topoi}, in section \ref{section: topoi}. A generalization of topoi, called \emph{(infinitary) pretopoi}, are introduced in \ref{section: pretopoi}. Afterwards, in section \ref{section: condensed sets}, we will start to develop the theory of condensed sets and how they relate to traditional topological spaces.

\section{Sites \& Sheaves}
\label{section: sites and sheaves}

A site is a generalization of the notion of open sets and their covers in a topological space. A site provides the minimal amount of data needed to construct a theory of sheaves. Open sets and their inclusions are modeled by a category, while covers of objects are specified as extra datum.

Much of this section follows Johnstone's \emph{Sketches of an Elephant, Volume II} \cite[C2]{elephant2}.


\begin{definition}[Coverages \& Sites, {
	\cite[Definition C2.1.1]{elephant2}}]
	\label{definition: coverages and sites}
	
	Let $\site C$ be a category. A \emph{coverage} on $\site C$ assigns to each $U \in \site C$ a collection $\covers U$ of families $(f_i\colon U_i \to U)_{i\in I}$ over $U$, called \emph{covers}, such that:
	
	If $(f_i\colon U_i \to U)_{i\in I}$ is a cover and $g\colon V\to U$ is any morphism in $\site C$, then there exists a cover $(h_j\colon V_j \to V)_{j\in J}$ of $V$, such that for each $j\in J$ the morphism $g\circ h_j$ factors through some $f_i$ with $i \in I$. That is every solid diagram 
	\begin{center}
		\begin{tikzcd}
			V_j\ar[r, dotted]\ar[d, "h_j"] &U_i\ar[d, "f_i", dotted]\\
			V\ar[r, "g"] &U
		\end{tikzcd}
	\end{center}
	can be completed by a pair of dotted arrows for some $i\in I$.
	
	A category equipped with a coverage is called a \emph{site}. A \emph{(locally) small site} is a site whose underlying category is (locally) small.
\end{definition}

\begin{remark}[Terminology]
	A \emph{cover} (recouvrement) $(f_i\colon U_i \to U)_{i\in I}$ of some object $U$ should be thought of as an open cover of $U$ by the $U_i$. A \emph{coverage} is a collection of covers. These notions should be distinguished from the unrelated \emph{covering} (revêtement) as \eg the universal covering $\R \to S^1$ of the circle $S^1$ in topology. Differentiating between the terms \emph{cover}, \emph{coverage} and \emph{covering} has provided the author with some headache.
\end{remark}

\begin{remark}
	The existence condition in definition \ref{definition: coverages and sites} essentially asserts that a kind of base change is possible -- for every cover of $U$ and any morphism $g\colon V\to U$, there is a cover of $V$ compatible with $g$. Indeed, if $\site C$ admits fiber products such that pullbacks of covers are covers, then this condition is met immediately.
\end{remark}

\begin{example}[Examples of sites]
	We now give two important (classes of) examples of sites.
	\begin{enumerate}
		\item
		Every topological space $X$ with its associated category $\open X$ of open sets naturally admits the structure of a site. Indeed, for every open $U \in \open X$ call a family $(U_i \to U)_{i \in I}$ a cover of $U$ if and only if $U = \bigcup_{i \in I} U_i$. 
		
		\item	
		The \emph{(small) étale site} $\etalesite X$ of a scheme $X$: Denote by $\Schemes$ the category of schemes. As a category, $\etalesite X$ is the full subcategory of $\Overcategory \Schemes X$ of étale morphisms over $X$. A cover of such an étale morphism $U\to X$ is a family $(f_i \colon U_i \to U)_{i\in I}$ over $X$ such that each $f_i$ is étale and the $f_i$ are jointly surjective, \ie the images of the $f_i$ cover $U$.
	\end{enumerate}
\end{example}

It is often important to consider sub-sites, that is subcategories of the underlying category equipped with the induced coverage.
\begin{definition}[Sub-sites]
	Let $\site C$ be a site and $\site C'$ be a subcategory. Declare a family $(U_i\to U)_{i\in I}$ in $\site C'$ to be a cover if it is one in $\site C$. Then $\site C'$ equipped with this coverage is called a \emph{sub-site of $\site C$}. If $\site{C'}$ is a full subcategory, then $\site{C}'$ is called a \emph{full sub-site of $\site{C}$}.
\end{definition}

We can now define presheaves and sheaves on a site. Inspired by the example of a topological space and its category of open sets, a presheaf should map each object of the underlying site (\ie an open set) to an object of a certain category (\eg $\Set$) while morphisms (\ie inclusions of open sets) should be mapped to morphisms in a contravariant manner (\ie to \emph{restrictions} between these associated objects). Hence, presheaves should be contravariant functors on the site.
\begin{definition}[Presheaves]
	Consider a site $\site C$ and a category $\category D$. A \emph{$\category D$-valued presheaf on $\site C$} is a functor $\site C^\op \to \category D$. Morphisms of presheaves are natural transformations. Hence, the category $\Presheaves{\site C}{\category D}$ of $\category D$-valued presheaves on $\site C$ is simply the functor category $\functorCategory{\site C^\op}{\category D}$.
\end{definition}

A presheaf of sets is now called a sheaf if it respects covers. In essence if $(U_i\to U)_i$ is a cover then the values of the (pre)sheaf on $U$ should be precisely determined by \emph{compatible} values on each of the $U_i$.
\begin{definition}[Compatible families \& Sheaves of sets, {
		\cite[Definition C2.1.2]{elephant2}}]
	\label{definition: compatible families and sheaves of sets}
	
	Consider a site $\site C$.
	
	Suppose $(f_i\colon U_i\to U)_{i\in I}$ is a cover in $\site C$ and $\sheaf F$ is a presheaf of sets on $\site C$. A family of sections $s_i\in \sheaf F(U_i)$ is called \emph{compatible} if for every pair of maps $U_i \xleftarrow{g} V \xrightarrow{h} U_j$ with $f_i\circ g = f_j\circ h$ in $\site C$ we have $\sheaf F(g)(s_i) = \sheaf F(h)(s_j)$, \ie the restrictions of $s_i$ and $s_j$ to $\sheaf F(V)$ agree.
	
	The presheaf of sets $\sheaf F$ is called a \emph{sheaf} on $\site C$ if for every cover $(f_i\colon U_i\to U)_{i\in I}$ and every family $(s_i)_{i\in I}$ compatible with the cover, there exists a unique $s\in \sheaf F(U)$ such that $\sheaf F(f_i)(s)=s_i$, \ie the family of sections glues to a unique section over $U$. We obtain the full subcategory
	$$\Sheaves{\site C}{\Set} \subseteq \Presheaves{\site C}{\Set}$$
	of sheaves of sets on $\site C$.
\end{definition}

We will also consider sheaves with values in categories other than sets -- however only in the special case where values are sets with extra structure. What this should mean we shall formalize in the following definition.
\begin{definition}[Very concrete categories]
	\label{definition: very concrete categories}
	
	A category $\category{D}$ is called \emph{very concrete} if it admits all filtered colimits and there is a continuous, faithful and conservative functor $\category{D} \to \Set$ that commutes with filtered colimits.
\end{definition}

\begin{remark}
	Definition \ref{definition: very concrete categories} has two purposes: First, it allows for the general theory of sheaves of sets to translate well to sheaves with values in the very concrete category, see definition \ref{definition: sheaves in general}. Secondly, it provides a rigid setting such that remark \ref{remark: transitioning in general} will apply -- this is fundamental to work with \emph{condensed objects in the very concrete category} (we will see later what this means).
\end{remark}

\begin{example}[Some very concrete categories]
	Very concrete categories are for example the category of groups, its subcategory of abelian groups, the category of rings but not the category of topological spaces (\eg $[0, 2\pi) \to S^1$ is bijective but not a homeomorphism, so the forgetful functor is not conservative).
\end{example}

We obtain the following straightforward generalization of definition \ref{definition: compatible families and sheaves of sets}.
\begin{definition}[Sheaves in general, {
	\cite[Definition C2.1.2]{elephant2}}]
	\label{definition: sheaves in general}
	
	Consider a site $\site C$ and a very concrete category $\category D$. Then a $\category D$-valued presheaf $\sheaf F$ is called a \emph{sheaf}, if the associated underlying presheaf of sets is a sheaf. We obtain a full subcategory
		$$\Sheaves{\site C}{\category D} \subseteq \Presheaves{\site C}{\category D}$$
	of $\category D$-valued sheaves on $\site C$.
\end{definition}

If the site admits fiber products (which one should regard as the analogue of the intersection of open sets), then compatibility of sections is given by \emph{agreement on intersections} and one can write down certain equalizer conditions that are sufficient and necessary for being a sheaf.
\begin{lemma}[Being a sheaf is an equalizer condition]
	\label{lemma: being a sheaf is an equalizer condition}

	Let $\site C$ be a site with fiber products and $\category D$ be a very concrete category. A presheaf $\sheaf F$ is a sheaf if and only if for each $U \in \site C$ and each cover $(f_i \colon U_i \to U)_{i\in I}$ of $U$ the natural morphism
	$$\begin{tikzcd}
		\sheaf F(U) \ar[r] &\equalizer\Bigg(\prod_{i\in I} \sheaf F(U_i) \ar[r, shift left=1.1] \ar[r, shift right=1.1] &\prod_{i, j \in I} \sheaf F(U_i \times_U U_j)\Bigg)
	\end{tikzcd}$$
	is an isomorphism, where the two parallel maps are induced by the assignments $(s_i)_{i\in I}\mapsto \sheaf F(U_i\times_U U_j \to U_i)(s_i)$ and $(s_i)_{i\in I}\mapsto \sheaf F(U_i\times_U U_j \to U_j)(s_j)$.
\end{lemma}
\begin{proof}
	Let $(f_i\colon U_i\to U)_i$ be a cover. Clearly $\prod_{i\in I} \sheaf F(U_i)$ is the set of families of sections $(s_i)_{i\in I}$ where $s_i\in \sheaf F(U_i)$.
	
	Suppose such a family $(s_i)_{i\in I}$ is compatible with the cover. Then by definition of compatibility, for all $i,j\in I$ the restrictions of $s_i$ and $s_j$ under the two maps $\sheaf F(U_i)\rightrightarrows \sheaf F(U_i\times_U U_j)$ agree. Hence, every compatible family lies in the equalizer.
	
	Now suppose we are given a family $(s_i)_{i\in I}$ in this equalizer. We will show that the family is compatible. Suppose that $U_i\xleftarrow g V\xrightarrow h U_j$ is any pair of morphisms with $f_i\circ g=f_j\circ h$. Then by the universal property of the fiber product there is a unique compatible morphism $u\colon V\to U_i\times_U U_j$, such that the diagram
	$$\begin{tikzcd}
		V\ar[dr, "u"]\ar[drr, bend left=15, "g"]\ar[ddr, bend right=15, "h", swap]&&\\
			&U_i\times_U U_j \ar[r, "p_i", swap] \ar[d, "p_j"] &U_i\ar[d, "f_i", swap]\\
			&U_j \ar[r, "f_j"] & U
	\end{tikzcd}$$
	commutes. Now as by definition of the equalizer we have $\sheaf F(p_i)(s_i)=\sheaf F(p_j)(s_j)$ it follows, that
		$$\sheaf F(g)(s_i)=(\sheaf F(u)\circ \sheaf F(p_i))(s_i)=(\sheaf F(u)\circ \sheaf F(p_j))(s_j)=\sheaf F(h)(s_j)$$
	proving compatibility, since $i, j$ were arbitrary.
	
	Hence, the equalizer is precisely the set of compatible families. The claim now follows immediately. Indeed, surjectivity of the induced map to the equalizer ensures that gluing is possible, while injectivity ensures that gluing is unique.
\end{proof}

As we are used to, there exists a canonical way of turning a presheaf of sets into a sheaf of sets -- this is of course sheafification.
\begin{definition}[Sheafification]
	\label{definition: sheafification}
	
	For a small site $\site C$ the inclusion $\Sheaves{\site C}{\Set} \to \Presheaves{\site C}{\Set}$ of sheaves into presheaves admits a left adjoint
		$$\sheafification{-}\colon \Presheaves{\site C}{\Set} \to \Sheaves{\site C}{\Set}$$
	that commutes with finite limits, called \emph{sheafification}.
\end{definition}
\begin{proof}
	This follows from \cite[Prop. C2.2.6]{elephant2} after observing that the colimits in the definition of the \emph{separation} $\sheaf F\mapsto\sheaf F^+$ (and of the sheafification $\sheafification{-}\colon \sheaf{F}\mapsto \sheaf{F}^{++}$) are filtered and hence commute with finite limits.
\end{proof}

\begin{remark}[Complications of sheafification in a general setting]
	We will see more general versions of definition \ref{definition: sheafification} later on. One is often interested in sheaves with values in some category $\category D$ that is not $\Set$. This in general turns out to be somewhat problematic and one needs to impose some conditions on the category $\category D$. A sufficient (but rather technical, hence we will not untangle them) set of conditions on $\category D$ asks, that $\category D$ is bicomplete, that in $\category D$ small filtered colimits are exact and that $\category D$ satisfies the so-called \emph{IPC-property}.
	
	We will not need that kind of generality as we will only consider (pre)sheaves of sets, abelian groups and rings. It is not very hard to see that presheaves of abelian groups/rings are equivalent to abelian/ring objects internal to the category of presheaves of sets. Since sheafification preserves finite limits one naturally obtains the structure of an abelian group/ring on the sheafification of the underlying presheaf. We will repeat this argument once we introduce abelian sheaves in chapter \ref{chapter: rings, modules and condensed algebra}.
\end{remark}

We should now ask ourselves how to relate different sites, in particular how to relate sheaves thereon. Once again, we turn to topological spaces for our defining motivation. If $f\colon X \to Y$ is a map of topological spaces, then by taking preimages under $f$ we obtain a functor $f^{-1}\colon \open Y \to \open X$ mapping unions to unions and covers to covers. As we are dealing with the sites associated to $X$ and $Y$, we only have access to open sets, not elements thereof, so this map $f^{-1}$ alone should determine a map of sites. However, a site should be considered as a generalized space, so the functor $f^{-1}$ \quote{points in the wrong direction}, as we ideally would like any continuous map $f\colon X\to Y$ to have a corresponding map of sites $\open X\to \open Y$ in the same direction. One solves this problem rather pragmatically by declaring a map of sites $f\colon \site{C}\to\site{D}$ to be a functor $f^\top\colon\site{D}\to \site{C}$ in the other direction, satisfying some additional properties. Certainly $f^\top$ should preserve covers. Now in the case of topological spaces, $f^{-1}$ maps intersection to intersections and more generally limits in $\open{Y}$ to limits in $\open{X}$. We will require the same of $f^\top$ if(!) $\site{D}$ is finitely complete. If $\site{D}$ does not admit all finite limits, then one has to be more careful.
\begin{definition}[Maps of sites]
	\label{definition: maps of sites}
	
	Let $\site{C}$ and $\site{D}$ be sites. Suppose that $\site{D}$ is finitely complete. Then a map $f\colon \site C \to \site D$ of sites is a functor $f^\top\colon \site D \to \site C$, commuting with finite limits and mapping covers to covers.
\end{definition}

\begin{remark}[Maps of sites in full generality]
	\label{remark: maps of sites in full generality}
	
	If the site $\site{D}$ in definition \ref{definition: maps of sites} is not finitely complete, one needs an alternative definition of \emph{maps of sites}. A generalization is for example given in \cite[Def. 4.10]{shulman2012exactcompletionssmallsheaves} by Shulman and asks for $f^\top$ to be \emph{covering flat}, essentially requiring that $f^\top$ \emph{preserves all finite limits that should have existed}. We will however only need this generalization in case that $\site{D}$ is a full subcategory of $\site{C}$. In this case it follows by \cite[Def 11.1 \& Thm. 11.2 (b)]{shulman2012exactcompletionssmallsheaves} that the inclusion satisfies this condition, hence giving us a map of sites $\site{C} \to \site{D}$. For a discussion of the notion of covering flatness see  \href{https://golem.ph.utexas.edu/category/2011/06/flat_functors_and_morphisms_of.html}{Shulman's post in the $n$-café}.
\end{remark}

Now having introduced maps of sites we can of course ask to transport sheaves along them. As with topological spaces, if $f\colon X\to Y$ is a continuous map and $\sheaf F$ is a sheaf on $X$, we easily obtain a sheaf on $Y$ by simply setting $f_\ast\sheaf F(V) \coloneqq \sheaf F(f^{-1}(V))$, \ie by precomposing with $f^{-1}$. This works since taking preimages maps open sets to open sets and preserves inclusions. Since in the general case taking preimages of open sets is modeled by the functor $f^\top$, it is clear how to proceed.
\begin{lemma}[Direct images]
	
	Let $f\colon \site C\to \site C'$ be a map of sites and $\category D$ a very concrete category. The functor
	\begin{align*}
		f_\ast \colon \Presheaves{\site C}{\category D} &\to \Presheaves{\site C'}{\category D},\\
		\sheaf F &\mapsto \sheaf F\circ f^\top \ldef \left((\site C')^\op \xrightarrow{f^\top}\site C^\op\xrightarrow{\sheaf F} \category D\right)
	\end{align*}
	is called the \emph{direct image} and restricts to a functor $f_\ast \colon \Sheaves{\site C}{\category D} \to \Sheaves{\site C'}{\category D}$.
\end{lemma}
\begin{shortproof}
	This is very similar to \cite[Lem. C2.3.3]{elephant2} keeping in mind \cite[Rem. C2.3.7]{elephant2}.
\end{shortproof}

Similarly, one can ask to transport sheaves in the other direction. As with topological spaces, this is initially problematic. If one tries to mimic the definition of the direct image by setting $f^{-1}\sheaf F(U) \coloneqq \sheaf F(f(U))$ one immediately runs into the problem that $f(U)$ might not be open and hence $\sheaf F(f(U))$ is ill-defined. Abstractly the solution comes from simply requesting a functor $f^{-1} \ladj f_\ast$ left adjoint to the direct image. More concretely, one can try to approximate $f(U)$ from above by open sets. For this take open sets $V$ such that $U\subseteq f^{-1}(V)$ (which implies $f(U)\subseteq V\cap \image f \subseteq V$) and glue the values $\sheaf F(V)$ suitably. In the general case this is described by a filtered colimit given by the comma-category of maps $U\to f^\top(V)$.
\begin{proposition}[Inverse images]
	\label{Proposition: Existence of inverse image}
	
	Let $f\colon \site C\to \site C'$ be a map of small sites and $\category D$ a very concrete category such that sheafification of $\category D$-valued presheaves exists. Assume that $\site C$ is small and that $\site C'$ is essentially small. There exists a functor $f^{-1}\colon \Sheaves{\site C'}{\category D}\to \Sheaves{\site C}{\category D}$ left adjoint to the direct image $f_\ast$, called the \emph{inverse image}. 
	
	More precisely for each $U \in \site C$ define the category $I_U$ with objects $\{(V, \phi)\setseparator V\in\site C', \phi\colon U \to f^\top(V)\}$ and morphism $\Hom_{I_U}((V, \phi), (V', \phi')) = \{g\colon V\to V'\setseparator g\circ \phi = \phi'\}$. For a sheaf $\sheaf F$ on $\site C'$, we obtain the inverse image as the sheafification
		$$f^{-1}\sheaf F = \sheafification{U\mapsto \colimit_{(V, \phi) \in I_U} \sheaf F(V)}.$$
	Notice how each morphism $g\colon U\to U'$ in $\site C$ functorially induces a morphism $I_U \to I_{U'}, (V, \phi)\mapsto (V, \phi\circ g)$, so the rule above similarly extends to morphisms, hence prior to sheafification we really obtain a presheaf.
\end{proposition}
\begin{shortproof}
	This is very similar to \cite[Cor. C2.3.4]{elephant2} keeping in mind \cite[Rem. C2.3.7]{elephant2}.
\end{shortproof}

In the special case where one site is a full sub-site of the other (see remark \ref{remark: maps of sites in full generality}), the situation is a bit more straightforward. For topological spaces this corresponds to pulling back from a coarser topology to a finer one, along the identity. In this case we cover $f(U)=U$ by open sets $V$ in the coarser topology.
\begin{remark}[Inverse images for inclusions]
	\label{remark: inverse images for inclusions}
	Suppose $f\colon\site C \to \site C'$ is the map of sites corresponding to a full inclusion $\site C'\subseteq \site C$ of underlying categories. Then the categories $I_U$ considered in proposition \ref{Proposition: Existence of inverse image} drastically simplify and we are justified in writing
		$$f^{-1}\sheaf F = \sheafification{U\mapsto \colimit_{\substack{U\to V\\V\in \site C'}} \sheaf F(V)}$$
	for any sheaf $\sheaf F$.
\end{remark}


\section{Topoi}
\label{section: topoi}

We will now study the properties of sheaf categories of small sites.
\begin{definition}[Grothendieck Topoi, {\cite[Definition C2.2.9]{elephant2}}]
	A \emph{(Grothendieck) topos} is a category $\category T$ equivalent to the category of sheaves $\Sheaves{\site C}{\Set}$ on some small site $\site C$.
\end{definition}

\begin{remark}[Being a topos is a property]
	\label{remark: giraud's theorem}
	
	In \ref{theorem: girauds axioms} we will see a theorem due to Giraud that characterizes a topos $\category T$ by certain axioms internal to the category $\category T$. Hence, \emph{being a topos} is a property of a given category, not extra structure. One of these axioms (or rather a relaxation thereof) will play an important role in the theory of condensed mathematics. The axiom states: $\category T$ admits a generating set. We will introduce the notion of generation in \ref{section: generation} and use it to provide the mentioned relaxation in \ref{section: pretopoi}.
\end{remark}

We will now define a suitable notion of morphism between topoi. Recall that for any map of sites $f\colon \site C\to\site D$ there is an adjoint pair $f^{-1}\ladj f_\ast$ between the associated topoi. We mimic such pairs and their properties when defining morphisms of topoi.
\begin{definition}[Morphisms of topoi, {\cite[Definition A4.1.1]{elephant1}}]
	\label{definition: Morphisms of topoi}
	
	A \emph{(geometric) morphism} of topoi $\category T\to \category T'$ is a pair $f^{-1} \ladj f_\ast$ of adjoint functors such that $f^{-1}\colon \category{T'}\to\category{T}$ commutes with finite limits. In particular, isomorphisms of topoi are precisely adjoint equivalences. We will denote such a morphism by
		$$\category T \xleftrightarrows{f^{-1}}{f_\ast} \category T'.$$
\end{definition}

\begin{remark}
	Due to the adjunction $f^{-1} \ladj f_\ast$ in definition \ref{definition: Morphisms of topoi}, we additionally know that $f^{-1}$ and $f_\ast$ commute with all colimits and all limits respectively. Hence, in total $f^{-1}$ is cocontinuous and finitely continuous, while $f_\ast$ is continuous.
\end{remark}

Now one can ask when a given sub-site gives rise to an equivalent category of sheaves. The analogue for topological spaces are of course bases. Recall that that sheaves on a topological space $X$ are the same as sheaves on a base of $X$ -- this is for example used to construct the structure sheaf of an affine scheme $\Spec A$ from its basis of distinguished opens (where one can explicitly give the sections on $D(f)$ with $f\in A$ as the localization $A_f$). We will restrict out attention to \emph{full sub-sites}, \ie sub-sites whose underlying category is full in the underlying category of the larger site.
\begin{definition}[Dense sub-site]
	Let $\site C$ be a site. A sub-site $\site B$ of $\site C$ is said to be \emph{dense}, if it is full and every $U\in \site C$ admits a cover $(f_i\colon U_i\to U)_{i\in I}$ where $U_i\in \site B$ for all $i\in I$.
\end{definition}

We now obtain the desired result which formalizes the idea, that giving a sheaf is the same as giving a sheaf on a base of that topology.
\begin{lemma}[Comparison lemma]
	\label{lemma: comparison lemma}
	Let $\site B$ be a small dense sub-site of a locally small site $\site C$. The map $\site C\to\site B$ of sites (given by the inclusion of the underlying categories) gives rise to an adjoint equivalence
		$$\Sheaves{\site C}{\Set}\xleftrightarrows{i^{-1}}{i_\ast} \Sheaves{\site B}{\Set}$$
	of topoi.
\end{lemma}
\begin{shortproof}
	This is \cite[C2.2.3]{elephant2} or alternatively \cite[Thm. 11.8]{shulman2012exactcompletionssmallsheaves} with suitable choice of $\kappa$.
\end{shortproof}

This lets us introduce some handy terminology. If the underlying category of a site is large, but the site has a small dense sub-site, we obtain from the comparison lemma \ref{lemma: comparison lemma} that their sheaf categories are equivalent. 
\begin{definition}[Essentially small sites]
	\label{definition: essentially small sites}
	
	A locally small site $\site C$ is called \emph{essentially small} if it admits a small dense sub-site.
\end{definition}

\begin{remark}[Terminological clash]
	Notice that there is some clash of terminology when talking about essentially small sites. When referring to an essentially small $\site C$, we might (contrary to the previous definition in definition \ref{definition: essentially small sites}) interpret this as saying that the underlying category of $\site C$ is essentially small (\ie its isomorphism classes are in bijection with a set). Indeed, some authors define an essentially small site to be a site whose underlying category is essentially small.
	
	Clearly, if the underlying category is essentially small, one can obtain a dense sub-site of $\site C$ by considering a skeleton of the underlying category with induced covers. This is then a small dense sub-site, so the site $\site C$ is essentially small. It is however not clear to the author that the converse holds true as well. To him there seems to be no reason that admitting a small dense sub-site tells anything about the isomorphism classes of the underlying category.
\end{remark}


\section{Generation}
\label{section: generation}

As mentioned in remark \ref{remark: giraud's theorem}, we will need to study the notion of generation. When defining condensed sets, we will need to work with sheaves on a large site. On such a site however, it is dangerous to consider the full category of sheaves, since there is no guarantee that it is locally small: For two given sheaves, morphisms between them are natural transformations -- these are collections of morphisms (the components of the transformation) indexed by the objects of the site in question! This is a priori a large amount of data! Indeed, the full category of sheaves on a large site in general is neither a topos nor locally small!

If one imposes certain regularity on the two sheaves, then one can biject the morphisms between any two regular sheaves to a set. Hence, the full subcategory of such regular sheaves turns out to be locally small a posteriori. Condensed sets indeed will be defined as regular sheaves on a large site! This regularity is what we will call \emph{smallness}, see definition \ref{definition: small (pre)sheaves}.

The question now remains: What are the properties of this new category of \emph{regular/small} sheaves? It turns out that for the category of condensed sets the answer is rather satisfactory! It satisfies all but one of Giraud's axioms: It does not admit a set of generators, but instead one obtains only a proper class of them! Such a category will be called a \emph{pretopos} -- we will introduce pretopoi in \ref{section: pretopoi}. 

Let us now work towards the notion of generation. A natural starting point is the theory of stalks. Recall that for any point $x$ of a topological space $X$, one obtains the stalk functor $(-)_x\colon \Sheaves{X}\Set\to \Set$. Let us model this in the setting of topoi. The inclusion $\ast \to X$ of the point induces a morphism of sites $i\colon\open{\ast} \to \open{X}$ which in turn induces a morphism of topoi $$\Set \iso \Sheaves{\ast}{\Set} \xleftrightarrows{i^{-1}}{i_\ast} \Sheaves{X}{\Set}.$$
Here $i^{-1}\sheaf F = \sheaf F_x$ for any sheaf $\sheaf F$ recovers the stalk of $\sheaf F$ at $x$. Such a geometric morphism will be called a \emph{point} of the topos. Recall that a morphism $f\colon \sheaf F\to\sheaf F'$ is an isomorphism if and only if $f_x\colon \sheaf F_x\to\sheaf F'_x$ is an isomorphism for all points $x\in X$. While in a general topos this will not be true, it is still instructive to consider points and stalks for a general topos nonetheless.
\begin{definition}[(Having enough) points in a topos]
	\label{definition: points in a topos}
	
	Let $\category T$ be a topos. A \emph{point} $x=(x^{-1}, x_\ast)$ of $\category T$ is a geometric morphism
		$$\Set\iso \Sheaves \ast \Set \xleftrightarrows{x^{-1}}{x_\ast}\category T.$$
	For $\sheaf F\in \category T$ the object $\sheaf F_x \ldef x^{-1}\sheaf F$ is called the \emph{stalk of $\sheaf F$ at $x$}. The topos $\category T$ is said to \emph{have enough points} if any morphism $f\colon \sheaf F\to \sheaf F'$ in $\category T$ is an isomorphism if and only if the morphism $f_x\ldef x^{-1}f$ of stalks is an isomorphism for each point $x$ of $\category T$.
\end{definition}

Unfortunately not all topoi have enough points. An example from logic is provided by Sina Hazratpour \href{https://sinhp.github.io/scribbling/22-11-2017-a-topos-without-points}{on his website}. However, even if they do, those points might be inexplicit and hard to work with. Furthermore, if one leaves the realm of topoi and works with sheaves on a large site instead, then it seems (see \cite{pointsPretopoiLargeSite}) that defining \emph{points} as in definition \ref{definition: points in a topos} is no longer appropriate.

We will thus abandon the notion of points and search for a suitable replacement: It turns out that the idea of considering functors that \emph{jointly} give information about the category in question is a fruitful approach on even a general category. For this we will ask for a collection of objects $G$ such that the functors $\Hom(G, -)$ are \emph{jointly faithful}. Such objects are called \emph{generators} or \emph{separators}. Clausen and Scholze make extensive use of generators in \cite{condenseddotpdf}: We will for example see in \ref{proposition: abelian sheaves form an abelian category} that the category of condensed abelian groups is abelian and has enough projectives -- the projectives being generators!

Now for a topos (\ie sheaves on a \emph{small} site) the collection of representable sheaves forms (essentially by Yoneda's lemma) a generating set. In particular one usually ask for a \emph{small} collection of generators. This however is problematic for applications in condensed mathematics (or more generally sheaves of a large site). We will thus need a generalization to possibly non-small collections. Unfortunately, Scholze and Clausen mention no such generalization in \eg \cite{condenseddotpdf} and the author is also not aware of any other standard reference. Instead, the author would like to thank \href{https://zll22.user.srcf.net/profile/}{Zhen Lin Low} for providing such a generalization in \cite{generationZhenLin}.
\begin{definition}[Generation]
	\label{definition: generation}
	
	Let $\category C$ be a category. A class $\mathcal G$ of objects of $\category C$ is called \emph{generating}, if for any object $X\in \category C$ there is a small(!) collection (\ie a set) $\mathcal G'\subseteq \mathcal G$, such that for any parallel pair of morphisms $f,g \colon X\to Y$ in $\category C$ with the property that $f\circ h = g\circ h$ for any map $h\colon G\to X$ from some $G\in\mathcal G'$, we already have $f=g$. Such a set $\mathcal G'$ is called \emph{separating for $X$}.
\end{definition}

\begin{remark}
	Suppose that $\category C$ is a category and $\mathcal G$ a generating class of $\category C$. Much like stalks, this generating class can make statements about a morphism $f\colon X\to Y$ being an isomorphism. Suppose we have a candidate morphism $g\colon Y\to X$ that we suspect is an inverse of $f$. This can be reformulated as asking the two questions $\id_X \overset{?}= g\circ f$ and $\id_Y\overset ? = f\circ g$, which can be answered using the generating class $\mathcal G$. Indeed, by checking $(\id_X)_\ast\overset ? = (g\circ f)_\ast \colon \Hom(G, X)\to \Hom(G, X)$ and $(\id_Y)\overset ? = (f\circ g)_\ast\colon \Hom(G, Y)\to\Hom(G, Y)$ for any $G\in\mathcal G'$ in some suitable separating set $\mathcal G'\subseteq \mathcal G$, we can answer these two questions!
\end{remark}

In case the category in question carries more structure, definition \ref{definition: generation} can be reformulated in a manner of ways. Of particular importance to us is the case of an abelian category.
\begin{lemma}[Generating classes in abelian categories]
	\label{lemma: generating classes in abelian categories}
	
	% This is from: Compare to https://mathoverflow.net/questions/377324/compact-object-and-compact-generator-in-a-category
	Let $\category A$ be an abelian category with arbitrary direct sums. Then for a class $\mathcal G$ of objects of $\category A$ the following are all equivalent:
	\begin{enumerate}[(i)]
		\item
		\label{generating sets in abelian categories: i}
		$\mathcal G$ is a generating class.
		
		\item
		\label{generating sets in abelian categories: epi}
		Every object $X\in\category A$ admits an epimorphism $\bigoplus_{i\in I} G_i \twoheadrightarrow X$ where $(G_i)_{i\in I}$ is a small(!) family of objects of $\mathcal G$.
		
		\item
		\label{generating sets in abelian categories: presentable}
		Every object can be realized as a coequalizer of a pair of morphisms between small direct sums of objects in $\mathcal G$.
	\end{enumerate}
\end{lemma}
\begin{proof}
	\ref{generating sets in abelian categories: i} $\Rightarrow$ \ref{generating sets in abelian categories: epi}:\newline
	Let $\mathcal G'\subseteq \mathcal G$ be a small separating set for $X$. Consider the morphism $p\colon \bigoplus_{\substack{G\to X\\G\in\mathcal G'}} G \to X$. We claim that this is an epimorphism. Indeed, if $f,g\colon X\to Y$ is a pair of parallel morphism in $\category A$ for which $f\circ p = g\circ p$, then also $f\circ h = g\circ h$ for any $h\colon G\to X$ with $G\in\mathcal G'$ since $h$ admits a factorization $G\to \bigoplus_{\substack{G\to X\\G\in\mathcal G'}} G \to X$ by definition of the direct sum. Thus by choice of $\mathcal G'$, we must have $f=g$. Since $f$ and $g$ were arbitrary, $p$ must be epic.
	
	\ref{generating sets in abelian categories: epi} $\Rightarrow$ \ref{generating sets in abelian categories: i}:\newline
	Let $X$ be any object in $\category A$. Choose an epimorphism $p\colon \bigoplus_{i\in I} G_i \twoheadrightarrow X$ and define $\mathcal G'\ldef \{G_i\setseparator i\in I\}$. We claim that the set $\mathcal G'$ is separating for $X$. For this let $f,g\colon X\to Y$ be a pair of parallel morphisms in $\category A$ such that for any $h\colon G\to X$ with $G\in\mathcal G'$ we have $f\circ h=g\circ h$. Then in particular for $i\in I$ with associated $h_i\colon G_i \to \bigoplus_{i\in I} G_i \twoheadrightarrow X$ we have $f\circ h_i = g\circ h_i$, so also $f\circ p = g\circ p$. But $p$ is epic, hence $f=g$.
	
	\ref{generating sets in abelian categories: epi} $\Rightarrow$ \ref{generating sets in abelian categories: presentable}:\newline
	Consider an epimorphism $p\colon \bigoplus_{i\in I} G_i \twoheadrightarrow X$. Pick a similar epimorphism $p'\colon \bigoplus_{j\in J} G_j' \twoheadrightarrow \ker p$ and define $q=i\circ p'$ where $i$ is the inclusion of the kernel. We obtain the commutative diagram
	$$\begin{tikzcd}
		0 &\bigoplus_{j\in J} G_j' 	&\bigoplus_{i\in I} G_i &\coKernel q 	&0\\
		0 &\kernel p 				&\bigoplus_{i\in I} G_i &X 				&0
		\ar[from=2-1, to=2-2]
		\ar[from=2-2, to=2-3, "i"]
		\ar[from=2-3, to=2-4, "p", two heads]
		\ar[from=2-4, to=2-5]
		\ar[from=1-1, to=1-2]
		\ar[from=1-2, to=1-3, "q"]
		\ar[from=1-3, to=1-4]
		\ar[from=1-4, to=1-5]
		\ar[from=1-2, to=2-2, two heads, "p'"]
		\ar[from=1-3, to=2-3, equal]
		\ar[from=1-4, to=2-4, dashed]
	\end{tikzcd}$$
	with exact rows and induced morphism $\coKernel q\to X$ by the universal property of $\coKernel q$ since $p \circ q = 0$. Hence, by the 5-lemma (with the induced morphism in the center), this induced morphism is an isomorphism, \ie $X$ is a cokernel of $q$.
	
	\ref{generating sets in abelian categories: presentable} $\Rightarrow$ \ref{generating sets in abelian categories: epi}:\newline
	Clear since any coequalizer is just a cokernel.
\end{proof}


\section{Pretopoi}
\label{section: pretopoi}

In this section we will provide the definition of pretopoi, where \emph{being a pretopoi} is a certain relaxation to \emph{being a topos}. As previously mentioned this relaxation hinges on a theorem of Giraud, stated in theorem \ref{theorem: girauds axioms}, which characterizes a topos by properties internal to the category in question. 

We will need a few additional concepts to state Giraud's theorem. These properties are true in any topos and are heavily inspired by the behavior of $\Set$. 

\begin{definition}[Universally disjoint coproducts]
	Let $\category T$ be a category with all finite limits and all small colimits.
	
	A coproduct $X\ldef \coprod_{i\in I} X_i$ in $\category T$ is called \emph{disjoint}, if the associated morphisms $X_i\to X$ are monomorphisms and for each pair $i\ne j$ in $I$ the fiber product $X_i\times_X X_j$ (the intersection of $X_i$ and $X_j$ inside $X$) is an initial object in $\category T$ (\ie the intersection is empty). If additionally for every morphism $X'\to X$ the pullbacks of the morphisms $X_i\to X$ induce an isomorphism
		$$\coprod_{i \in I} (X'\times_X X_i) \to X',$$
	then the coproduct is called \emph{universally disjoint}.
\end{definition}

\begin{definition}[Equivalence relations \& quotients]
	Let $\category T$ be a category with all finite limits and all small colimits.
	
	Let $X$ be an object of $\category T$. A \emph{relation on $X$} is a monomorphism $\iota\colon R\to X\times X$ from some object $R$ of $\category T$. Denote its components by $p_1, p_2\colon R\to X$, so that $\iota=(p_1, p_2)$. Then $\iota$ is called an \emph{equivalence relation}
	\begin{enumerate}[(i)]
		\item
		if the diagonal map $X\to X\times X$ factors over $\iota$ (that is $\iota$ is \emph{reflexive}),
		
		\item
		if the morphism $(p_2, p_1)\colon R\to X\times X$ factors over $\iota$ (that is $\iota$ is \emph{symmetric}),
		
		\item 
		and if the morphism $(p_1\circ q_1, p_2\circ q_2)\colon R\times_X R \to X\times X$ factors over $\iota$ as well, where $R\times_X R$, $q_1$ and $q_2$ are given by the pullback
			$$\begin{tikzcd}
				R\times_X R\ar[r, "q_2"]\ar[d, "q_1"] &R \ar[d, "p_1"]\\
				R \ar[r, "p_2"] &X
			\end{tikzcd}$$
		(that is $\iota$ is \emph{transitive}).
	\end{enumerate}
	If $\iota\colon R\to X\times X$ is an equivalence relation then its \emph{quotient} is the coequalizer
		$$R\rightrightarrows X\twoheadrightarrow E/R,$$
	and it is always an epimorphism.
\end{definition}

\begin{definition}[Kernel pairs]
	Let $\category T$ be a category with all finite limits and all small colimits.
	
	If $u\colon X\to X'$ is any morphism, then a \emph{kernel pair} for $u$ is a parallel pair of morphisms $p_1, p_2\colon R\to X$ such that
		$$\begin{tikzcd}
			R\ar[r, "p_2"]\ar[d, "p_1"] &X\ar[d, "u"]\\
			X\ar[r, "u"] &X'
		\end{tikzcd}$$
	is a pullback square.
\end{definition}

\begin{definition}[Stably exact forks]
	Let $\category T$ be a category with all finite limits and all small colimits.
	
	A diagram in $\category{T}$ of the form
		$$\begin{tikzcd}
			R\ar[r, "p_1", shift left=.25em]\ar[r, "p_2", shift right=.25em, swap] &X \ar[r, "q"] &Q
		\end{tikzcd}$$
	is called a \emph{fork}. A fork is said to be \emph{exact} if $q$ is a coequalizer of $p_1$, $p_2$ and if $p_1$, $p_2$ form a kernel pair for $q$. If additionally any base change
		$$\begin{tikzcd}
			R\times_Q Q' \ar[r, shift left=.25em]\ar[r, shift right=.25em] &X\times_Q Q' \ar[r] &Q \times_Q Q'\cong Q'
		\end{tikzcd}$$
	along some morphism $Q'\to Q$ stays exact, then the fork is called \emph{stably exact}.
\end{definition}

We can now state Giraud's axioms.
\begin{theorem}[Giraud's axioms]
	\label{theorem: girauds axioms}
	
	A locally small category $\category T$ with all finite limits and all small colimits is a topos if and only if
	\begin{enumerate}[(i)]
		\item
		coproducts in $\category T$ are universally disjoint,
		
		\item
		every epimorphism in $\category T$ is a coequalizer,
		
		\item
		every equivalence relation in $\category T$ is a kernel pair,
		
		\item
		every exact fork $R\rightrightarrows X\to Q$ in $\category T$ is stably exact,
		
		\item
		and $\category T$ admits a small generating class.
	\end{enumerate}
\end{theorem}
\begin{shortproof}
	This is proven in \cite[Appendix §3]{sheaves-in-geometry-and-logic}.
\end{shortproof}


This theorem now allows us to define what a \emph{pretopos} should be. As previously mentioned the idea is to relax Giraud's axioms by allowing \emph{large} classes of generators.
\begin{definition}[(Infinitary) pretopoi]
	\label{definition: pretopoi}
	
	A locally small category $\category T$ with all finite limits and all small colimits is called a \emph{(infinitary) pretopos} if
	\begin{enumerate}[(i)]
		\item
		\label{definition: pretopoi, coproducts are universally disjoint}
		coproducts in $\category T$ are universally disjoint,
		
		\item
		\label{definition: pretopoi, every epimorphism is a coequalizer}
		every epimorphism in $\category T$ is a coequalizer,
		
		\item
		\label{definition: pretopoi, every equivalence relation is a kernel pair}
		every equivalence relation in $\category T$ is a kernel pair,
		
		\item
		\label{definition: pretopoi, every exact fork is stably exact}
		every exact fork $R\rightrightarrows X\to Q$ in $\category T$ is stably exact,
		
		\item
		\label{definition: pretopoi, existence of a generating class}
		and $\category T$ admits a generating class.
	\end{enumerate}
	
	A \emph{(geometric) morphism of pretopoi} $\category T\to\category T'$ is a pair $f^{-1}\ladj f_\ast$ of adjoint functors, such that $f^{-1}$ commutes with finite limits. As with morphisms of topoi, we will denote such a morphism of pretopoi by
		$$\category T \xleftrightarrows{f^{-1}}{f_\ast} \category T'.$$
\end{definition}

\section{Regular Projectivity \& Compactness}

The category of condensed sets will be a pretopos, so admits a class of generators. This class of generators can be chosen to consist of very well-behaved objects. In this section we will introduce the terminology that formalizes what this is supposed to mean.
\begin{definition}[Reflexive coequalizers]
	A \emph{reflexive pair} in a category $\category T$ is a diagram of the form
		$$\begin{tikzcd}
			\cdot \ar[r, shift left=0.35em]\ar[r, shift right=0.35em] &\cdot\ar[l]
		\end{tikzcd}$$
	where the arrow to the left is a common section of the two arrows to the right. A colimit of a reflexive pair is called a \emph{reflexive coequalizer}.
\end{definition}

\begin{remark}[Reflexive pairs from reflexive relations]
	\label{remark: reflexive pairs from reflexive relations}
	
	Notice that if $X$ is some object and $\iota\colon R\to X\times X$ is any reflexive relation on $X$ then by definition there is some morphism $X\to R$ providing a factorization of $X\to X\times X$ along $\iota$. This morphism is a common section of the two components $R\to X$ of $\iota$. Hence, to every reflexive relation we can naturally associate a reflexive pair. If the relation is even an equivalence relation then the associated quotient is a reflexive coequalizer for this associated reflexive pair.
\end{remark}

\begin{definition}[Regular projectivity \& compactness]
	\label{definition: regular projectivity and compactness}
	
	Let $\category T$ be a cocomplete category and $P\in \category T$. Then $P$ is called
	\begin{itemize}
		\item
		\emph{(reflexive) projective} if $\Hom_\category T(P, -)$ commutes with reflexive coequalizers,
		
		\item 
		\emph{compact} if $\Hom_\category T(P, -)$ commutes with filtered colimits.
		
		\item 
		\emph{compact projective} if it is both compact and reflexive projective.
	\end{itemize}
\end{definition}

\begin{remark}[Comparison of reflexive projectivity to normal projectivity]
	The definition of reflexive projectivity in definition \ref{definition: regular projectivity and compactness} is closely related to the usual notion of projectivity and provides a relaxation to it. Instead of requiring the lifting against all epimorphisms, being a reflexive projective requires only the lifting against reflexive epimorphisms, that is morphisms that arise as coequalizers of reflexive pairs! By remark \ref{remark: reflexive pairs from reflexive relations} any quotient of an equivalence relation fits that description. Hence, being a reflexive projective allows to lift against \emph{gluing maps}.
	
	If the ambient category is a pretopos, then every epimorphism $q\colon X\to Y$ is already the quotient of the equivalence relation $X\times_Y X\to X\times X$, so is a reflexive epimorphism. Thus, in a pretopos reflexive projectivity and projectivity agree! In particular, we may simply speak of \emph{projectivity} throughout this thesis.
\end{remark}
